% !TeX program = lualatex
% Compile twice: lualatex 72.tex
% Standalone research notebook; original section preserved.
% Research additions: 27 September 2026. No external TeX input or BibTeX file.
\documentclass[11pt,a4paper]{article}
\usepackage[margin=24mm,headheight=14pt]{geometry}
\usepackage{fontspec}
\setmainfont{Latin Modern Roman}
\setsansfont{Latin Modern Sans}
\setmonofont{Latin Modern Mono}
\usepackage{amsmath,amssymb,mathtools}
\usepackage{unicode-math}
\setmathfont{Latin Modern Math}
\usepackage{microtype}
\usepackage{enumitem,xurl,needspace}
\usepackage[dvipsnames]{xcolor}
\usepackage{hyperref}
\hypersetup{unicode=true,colorlinks=true,linkcolor=MidnightBlue,urlcolor=MidnightBlue,
 pdftitle={Research notebook - Problem 72},pdfauthor={Open Problems Math},bookmarksnumbered=true}
\usepackage{bookmark}
\usepackage{titlesec}
\titleformat{\section}{\Large\bfseries\raggedright}{\thesection}{0.7em}{}
\usepackage{fancyhdr}
\pagestyle{fancy}\fancyhf{}
\fancyhead[L]{\small\sffamily Open Problems Math | Research notebook}
\fancyhead[R]{\small\sffamily Problem 72 | 27 September 2026}
\fancyfoot[C]{\thepage}
\setlength{\parindent}{0pt}
\setlength{\parskip}{5pt plus 1pt}
\setlength{\emergencystretch}{3em}
\setcounter{tocdepth}{1}
\setcounter{secnumdepth}{1}
\setlist[itemize]{leftmargin=3em,itemsep=5pt,topsep=4pt,parsep=0pt}
\allowdisplaybreaks
\numberwithin{equation}{section}
\newcommand{\N}{\mathbb{N}}
\newcommand{\Z}{\mathbb{Z}}
\newcommand{\Q}{\mathbb{Q}}
\newcommand{\R}{\mathbb{R}}
\newcommand{\C}{\mathbb{C}}
\newcommand{\F}{\mathbb{F}}
\newcommand{\A}{\mathbb{A}}
\newcommand{\PP}{\mathbb P}
\newcommand{\E}{\mathbb{E}}
\newcommand{\eps}{\varepsilon}
\newcommand{\dd}{\,\mathrm d}
\newcommand{\sref}[2]{\hyperlink{source:#1:#2}{[S#2]}}
\newcommand{\eref}[2]{\hyperlink{extra:#1:#2}{[E#2]}}
\newif\ifshowcatalogue
\showcataloguetrue
\newcommand{\cataloguescope}[1]{\ifshowcatalogue
 \paragraph{Exact scope in the supplied catalogue.}
 \begin{quote}\small #1\end{quote}\fi}
\begin{document}
\setcounter{section}{71}
\section{\texorpdfstring{Termination of Flips Conjecture in the Minimal Model Program}{Termination of Flips Conjecture in the Minimal Model Program}}\label{72}
\subsection{Definitions and mathematical statement}
A log pair $(X,\Delta)$ has $X$ normal, $\Delta$ an effective boundary divisor, and $K_X+\Delta$ real or rational Cartier in the chosen MMP category. Log canonicity requires all discrepancies to be at least $-1$; klt requires them to be greater than $-1$. A flip replaces a small negative contraction by the small birational model on which the transformed adjoint divisor is relatively ample. The claim excludes any infinite chain
\[
 (X_0,\Delta_0)\dashrightarrow(X_1,\Delta_1)\dashrightarrow\cdots
\]
of the allowed adjoint flips starting from a projective complex log canonical pair. A theorem about one prescribed choice of flips is not automatically termination of every sequence.
\par\smallskip\textit{Formulation and concept references:} \sref{72}{1}.
\cataloguescope{Prove that there is no infinite sequence of flips in the minimal model program: for a projective log canonical (or klt) pair over the complex numbers, every sequence of (K\_X + Delta)-flips terminates after finitely many steps, in every dimension.}
\subsection{Short English statement}
Must every allowed sequence of flips in the complex minimal model program stop after finitely many steps?
% BEGIN RESEARCH_ADDITION ATTEMPT_72
\subsection{Research attempt}

\paragraph{Current frontier.}
Lazi\'c--Xie connect arbitrary termination to Nakayama--Zariski decomposition
and the existence of terminating sequences under their hypotheses; this
is not a theorem that every conceivable sequence already terminates
\eref{72}{1}. Moraga's 2025 theorem concerns pseudo-effective log canonical
fourfold flips, not arbitrary dimension \eref{72}{2}.
The distinction between a terminating MMP selected with scaling and
termination of every sequence is central. Below a discrete monotone
certificate is tested without assuming that any such certificate is
available for all pairs.

\paragraph{Attack route.}
Canonical rings appear promising because flips are designed by adjoint
linear systems, but their sections are largely invariant under small
birational changes. Discrepancies can improve instead, suggesting a finite
accounting of their total improvement. A third route would control changing
centres geometrically and reduce to a finite collection of valuations.
The attempt makes that third requirement explicit inside a discrepancy
certificate, then checks why boundedness, monotonicity and rational
coefficients separately do not suffice.

\paragraph{Attempt: a finite discrepancy ledger.}
Work first with a fixed rational-boundary pair and a sequence whose models
have a common function field. Let $v_1,\ldots,v_r$ be finitely many
normalized divisorial valuations of that field. Write $a_i(j)$ for their
log discrepancies at the $j$th model. Suppose that there are an integer
$q\geq1$, finite upper bounds $B_i$, and an initial index $j_0$ such that
\begin{equation}\label{eq72:ledger}
 a_i(j)\in q^{-1}\Z,\quad
 a_i(j_0)\leq a_i(j)\leq B_i,\quad
 a_i(j+1)\geq a_i(j),
\end{equation}
and every flip after $j_0$ strictly increases at least one $a_i$.
Then the number of those flips is at most
\begin{equation}\label{eq72:bound}
 \left\lfloor q\sum_{i=1}^r(B_i-a_i(j_0))\right\rfloor.
\end{equation}
Indeed the integer-valued quantity
$q\sum_i(a_i(j)-a_i(j_0))$ is nonnegative, is bounded above by the
expression before the floor, and increases by at least one at every step.
This proves termination under the certificate. It handles arbitrary choices
of the subsequent flips, provided the same hypotheses hold for that
sequence; it does not merely construct one favourable choice.

The geometric discrepancy monotonicity motivating this setup follows from
the negativity comparison for flips, as used in the source
\eref{72}{1}. The statement just proved does not import more from that
comparison than monotonicity. In particular, it does not assume that a
fixed finite valuation set detects every flip, that the discrepancies of
those valuations have upper bounds, or that the indices of all models
have a common denominator. Those are precisely the additional conditions
that an actual construction of \eqref{eq72:ledger} must establish.

One slightly more flexible variant allows the valuations to be introduced
in stages, but assigns each newly introduced valuation a nonnegative
integer budget and requires a finite total budget for the entire sequence.
If every flip spends at least one unit and no unit is reused, the same
counting proof terminates the sequence. This restatement does not evade
the difficulty: a proof of a finite total budget would be the missing
geometric finiteness theorem, not a formal consequence of the existence of
only finitely many boundary components on the initial model.

\paragraph{Attempt: why the canonical ring supplies no such clock.}
Let $X\dashrightarrow X^+$ be a small birational map of normal projective
varieties, and let $D,D^+$ be corresponding rational divisors. At every
common sufficiently divisible degree $m$, global sections can be viewed
inside the common rational function field as
\[
 H^0(X,\mathcal O_X(mD))=
 \{g:\operatorname{div}_X(g)+mD\geq0\}\cup\{0\}.
\]
Smallness identifies all prime divisors and their valuations, as well as
the coefficients of the strict transforms. Thus the effectivity conditions
are identical and
\begin{equation}\label{eq72:ring}
 H^0(X,mD)=H^0(X^+,mD^+)
\end{equation}
in the common field. For non-Cartier degrees one can use the corresponding
reflexive divisorial sheaves; the same codimension-one argument applies.
Normality is what allows those divisorial conditions to characterize the
sections. Hence the adjoint graded ring is preserved under a flip in this
sense, and its Hilbert function does not strictly decrease.

Finite generation of a ring can help organize particular models, but
\eqref{eq72:ring} shows why the naive proof ``each flip simplifies the
canonical ring'' has no numerical decrease to count. One needs extra
control of chambers, centres, or local models. Equality of the rings at
every step is fully compatible with a hypothetical infinite chain of
different birational presentations; the equality alone says nothing about
its length. This is an obstruction to the proposed invariant, not a
constructed infinite flip sequence.

\paragraph{Stress test: four distinct missing uniformities.}
A strictly increasing bounded rational sequence need not stop:
$a(j)=1-2^{-j}$ is a counterexample. Thus the denominator condition in
\eqref{eq72:ledger} cannot be replaced merely by rationality of every term.
Even a single initial Cartier index does not, without proof, bound the
indices on every later model. Real-boundary versions need a different
well-foundedness mechanism and are not covered by a lattice of rational
values.

Conversely, a uniform positive increment and bounded individual coordinates
do not suffice with infinitely many coordinates. The vectors with the
first $j$ coordinates equal to one and all later coordinates zero increase
strictly forever, while each coordinate has range $\{0,1\}$. Each step
could be witnessed by a fresh valuation. This disproves a purely formal
inference from ``every flip improves some discrepancy'' to a finite
sequence. It is not a valuation construction on an actual variety.

Nor can upper bounds be dropped for a finite ledger: the one-coordinate
sequence $a(j)=j$ has all the other properties and never stops. Finally,
strict improvement must occur in the chosen ledger at every step; an
infinite chain leaving those finitely many coordinates unchanged is not
excluded by their boundedness. These examples separate the hypotheses
rather than conflating them into a vague appeal to discreteness.

The geometric challenge is therefore to produce a finite detecting family
or finite total budget, with uniform value control, after any allowable
initial segment. Bounds on one type of singularity, pseudo-effectivity in
dimension four, or a prescribed scaling construction are useful partial
inputs but do not verify that package in all dimensions and for all lc
pairs. No assertion is made here that minimal log discrepancies are known
to satisfy all the semicontinuity or chain conditions such a strategy might
require. Using an unresolved abundance or singularity conjecture would be
an additional conditional input, not a completion.

\paragraph{Outcome.}
\textbf{Outcome: Conditional reduction.}

The attempt proves an explicit finite-budget termination criterion and
shows that the canonical-section ring is not a decreasing invariant for
small flips. Four elementary countermodels identify the distinct ways in
which a naive discrepancy argument can fail. They are tests of a proof
strategy, not counterexamples to termination.

What remains is to construct the geometric certificate for every allowable
sequence in the original scope, including an appropriate replacement for
rational discreteness in any real-boundary convention. Neither that
construction nor a new unconditional termination case is supplied. The
counting lemmas are elementary notebook deductions without a novelty claim.

% Keep the local bibliography together rather than leave a source fragment.
\needspace{170mm}
% END RESEARCH_ADDITION ATTEMPT_72
\subsection{Sources}
\begin{itemize}\raggedright
\item[\textnormal{[S1]}]\hypertarget{source:72:1}{} Vladimir Lazic, Zhixin Xie, "Nakayama-Zariski decomposition and the termination of flips," arXiv:2305.01752v4; received July 8, 2024, accepted June 25, 2025.
\url{https://master-math-fonda.imj-prg.fr/gt/mmp.pdf}.
{\small\textit{Catalogue formulation source.}}
\item[\textnormal{[S2]}]\hypertarget{source:72:2}{} MMP for Enriques pairs and singular Enriques varieties.
\url{https://jep.centre-mersenne.org/item/10.5802/jep.334.pdf}.
{\small\textit{Catalogue research/status reference; not a proof endorsement.}}
\item[\textnormal{[S3]}]\hypertarget{source:72:3}{} Nakayama-Zariski decomposition and the termination of flips.
\url{https://epiga.episciences.org/16998}.
{\small\textit{Catalogue research/status reference; not a proof endorsement.}}
\item[\textnormal{[S4]}]\hypertarget{source:72:4}{} On termination of flips and exceptionally non-canonical singularities.
\url{https://arxiv.org/abs/2209.13122}.
{\small\textit{Catalogue research/status reference; not a proof endorsement.}}
% BEGIN RESEARCH_ADDITION REFERENCES_72
\item[\textnormal{[E1]}]\hypertarget{extra:72:1}{} Vladimir Lazi\'c and
Zhixin Xie, \emph{Nakayama--Zariski decomposition and the termination of flips},
arXiv:2305.01752v4; \emph{\'Epijournal de G\'eom\'etrie Alg\'ebrique} 9 (2025),
article 24, introduction and hypotheses of the termination reductions.
\url{https://arxiv.org/html/2305.01752v4}.
\url{https://epiga.episciences.org/16998}.
{\small\textit{Inspected 27 September 2026; existence of one terminating
sequence and universal termination are not interchanged.}}
\item[\textnormal{[E2]}]\hypertarget{extra:72:2}{} Joaqu\'in Moraga,
\emph{Termination of pseudo-effective 4-fold flips},
Mathematische Zeitschrift 310 (2025), article 73.
\url{https://link.springer.com/article/10.1007/s00209-025-03766-y}.
{\small\textit{Publisher theorem scope inspected 27 September 2026;
pseudo-effective log canonical fourfolds, not the all-dimensional target.}}
% END RESEARCH_ADDITION REFERENCES_72
\end{itemize}

\end{document}
